\scanpage{5}
\mspage{\deleted{\unsure{3}}4}
\tikz[baseline=-.4ex,x=.75em,y=.75em]{\draw (.5,.5) circle (.48);\draw (.18,.48)--(.82,.48) (.4,.2)--(.65,.8) (.32,.76)--(.72,.23);} $M={}$ ensemble des \unsure{``}\uline{figures} $\subset\mathscr F$ \uline{élémentaires}'' ou \uline{multistrates}

\margin{\tikz[x=.6em,y=.6em]{\draw (0,1.7)--(0,2.85) .. controls (0,3.06) and (.19,3.12) .. (.50,3.12)--(2.55,3.18) .. controls (2.80,3.19) and (2.85,3.62) .. (2.88,3.9); \draw (0,1.7)--(1.20,1.72)--(2.88,1.53)--(3.02,0);}\enspace\uline{strates} d'une figure : les $X\in M$ tels que $X\leq F$\\ \uline{déployement} de $F$ : ens. de ses \uline{strates}\\ $\widetilde F=\{X\in M\mid X\leq F\}$\\ On écrit $X\triangleleft F$ ($X$ \uline{incident} à $F$) pour $X\in M$ et $X\leq F$}
\margin{formé des objets de $(\mathscr F,\leq)$ strictement irréductibles. Muni des relations $\underset{M}{\leq}$, $\underset{M}{\ll}$ induites par \deleted{\unsure{$\mathscr F$}} $\leq$ et $\ll$ dans $\mathscr F$}

\tikz[baseline=-.4ex,x=.75em,y=.75em]{\draw (.5,.5) ellipse (.43 and .55);\draw (.3,.15)--(.3,.8) (.2,.5)--(.8,.5) (.3,.66)--(.58,.9) (.3,.45)--(.62,.18);} $L\subset M$, ens. des \deleted{\uline{\unsure{éléments|points}}} \inserted{``lieux''} \deleted{\unsure{minimaux}}
\inserted{\tikz[x=.7em,y=.7em]{\draw (0,1.27)--(1.72,1.22) .. controls (2.06,1.22) and (2.42,1.02) .. (2.75,.85)--(4.67,.04)..controls(4.87,.17)and(5.17,.06)..(5.43,.08); \draw (2.57,.93)..controls(2.80,.54)and(4.13,.21)..(4.67,.04);}\enspace éléments minimaux de $\mathscr F\setminus\{\varnothing_{\mathscr F}\}$ pour $\ll$}

\uline{Axiomes}

I) Axiomes ``sine qua non''. Les premiers trois \inserted{At~1 - At~3} ne concernent que $(\mathscr F,\leq)$, les quatre suivants At~4 - At~7 \deleted{\unsure{At 1|At 4}} \deleted{\unsure{et|4}} et les relations entre $\leq$ et $\ll$

% slanted
\margin{Axiomes des \uline{sommes majorées de figures}}
\uline{At 1}\quad $\mathscr F\neq\varnothing$ et stable par $\underset{I}{\mathrm{Sup}}$ majorés ($I$ quelconque)

% slanted
\margin{Axiome des strates équivaut à :\\ Si $F,C\in\mathscr F$, alors\\ $F\leq G\Longleftrightarrow\forall X\in\widetilde F,\ \exists Y\in\widetilde G$ avec $X\leq Y$\\ Critère d'\uline{inclusion de figures par strates}}
\uline{At 2}\quad Toute figure \deleted{\unsure{d}}est le sup de ses strates :
\[
 F=\underset{X\in\widetilde F}{\mathrm{Sup}}\,X\qquad\text{où}\qquad
 \widetilde F=\{X\in M\mid X\leq F\}
\]

\uline{Scholie}\qquad \uline{At}~1 et \uline{At}~2 signifient que $\mathscr F$ s'interprète, via $F\mathrel{\tikz[baseline=-.45ex,x=.7em,y=.45em]{\draw (0,-.44)--(0,.44) (0,0)--(1.2,0) (1.2,-.44)--(1.2,.44);}}\widetilde F$, comme une partie \inserted{$\widetilde{\mathscr F}$} de $\mathscr F(M)$, $M$ ens ordonné (pour $\leq$), avec les conditions :

a) les $\Phi\in\widetilde{\mathscr F}$ sont fermés

b) Si $\Phi\in\widetilde{\mathscr F}$, toute partie fermée de $\Phi$ est $\in\widetilde{\mathscr F}$

c) $\deleted{\widetilde{\mathscr F}\neq\varnothing\quad(\Longleftrightarrow\varnothing\in\widetilde{\mathscr F})}$ (\deleted{\unsure{et}} $\displaystyle\bigcup_{\Phi\unsure{\in|+}\widetilde{\mathscr F}}\Phi=M$, et $\widetilde{\mathscr F}\neq\varnothing$
% slanted
\margin{$\widetilde{\mathscr F}\neq\varnothing$ automatique si $M\neq\varnothing$ et dans tous les cas, \unsure{équivaut} à $\varnothing\in\widetilde{\mathscr F}$}

ou, indépendamment d'une relation d'ordre sur $M$, un ens de parties satisfaisant
\diagram{\begin{tikzcd}[ampersand replacement=\&,column sep=0pt,row sep=2ex]
{} \& |[align=left,text width=.83\linewidth]| \text{b$'$)}\quad \forall x\in M,\ \text{l'ens } \widetilde{\mathscr F}^{x}=\{\deleted{\unsure{F|A}}\Phi\in\overline{\mathscr F}\mid x\in\Phi\}\ \text{a un} \\
{} \& |[align=left,text width=.83\linewidth]| \text{plus petit élément }\deleted{F}\ \text{(noté }\widetilde X\text{)}\\
{}\arrow[uu,bend left=35] \& |[align=left,text width=.83\linewidth]| \text{a$'$)}\quad\deleted{\unsure{\forall|F|\forall F}}\ \text{Dans }(\widetilde{\mathscr F},\subset),\ \text{toutes famille}\\
{} \& |[align=left,text width=.83\linewidth]| \text{majorée a un Sup, \uline{et} }\widetilde{\mathscr F}\neq\varnothing\\
{} \& |[align=left,text width=.83\linewidth]| (\Longleftrightarrow\varnothing\in\widetilde{\mathscr F})
\end{tikzcd}}

Relation
\[
\begin{aligned}
F\mathrel{\tikz[baseline=-.4ex,x=.65em,y=.65em]{\draw (0,1.1)--(1,.5)--(0,-.1) (1,1.1)--(0,.5)--(1,-.1);}}G
&\overset{\text{\unsure{dfn|défn|defn}}}{\Longleftrightarrow}
 F\vee C\ \text{existe i.e. }\{F,C\}\ \text{majoré dans }(\mathscr F,\leq)\\
&\text{\uline{figures} \uline{compatibles}}
\end{aligned}
\]
cette relation induit $X\mathrel{\tikz[baseline=-.4ex,x=.65em,y=.65em]{\draw (0,1.1)--(1,.5)--(0,-.1) (1,1.1)--(0,.5)--(1,-.1);}}Y$ dans $M$. Dans le cas où les figures sont de t.f., alors $(\mathscr F,\leq)$ \deleted{\unsure{est}} \inserted{\tikz[x=1em,y=1em]{\draw (8,1.2) .. controls (8,1) and (8,1) .. (7.7,1)--(.6,1) .. controls (.2,1) and (.2,.3) .. (0,.3);\draw (.45,.6)--(0,.3)--(.55,.28);\draw (.3,.7)--(0,.3)--(.45,.2);}\enspace avec At~1, At~2, At~3} se reconstitue par
\[
(M,\leq,\mathrel{\tikz[baseline=-.4ex,x=.65em,y=.65em]{\draw (0,1.1)--(1,.5)--(0,-.1) (1,1.1)--(0,.5)--(1,-.1);}}),\quad\text{où }\mathrel{\tikz[baseline=-.4ex,x=.65em,y=.65em]{\draw (0,1.1)--(1,.5)--(0,-.1) (1,1.1)--(0,.5)--(1,-.1);}}\text{ satisfait}
\]
\[
\left\{\begin{aligned}
&X\ \mathrel{\tikz[baseline=-.4ex,x=.65em,y=.65em]{\draw (0,1.1)--(1,.5)--(0,-.1) (1,1.1)--(0,.5)--(1,-.1);}}\ \text{réflexive symétrique}\\
&X\mathrel{\tikz[baseline=-.4ex,x=.65em,y=.65em]{\draw (0,1.1)--(1,.5)--(0,-.1) (1,1.1)--(0,.5)--(1,-.1);}}Y,\ X'\leq X,\ Y'\leq Y\ \Longrightarrow\ X'\mathrel{\tikz[baseline=-.4ex,x=.65em,y=.65em]{\draw (0,1.1)--(1,.5)--(0,-.1) (1,1.1)--(0,.5)--(1,-.1);}}Y'
\end{aligned}\right.
\]
